\documentclass[10pt]{letter}
\usepackage[letterpaper,margin=0.85in]{geometry}
\usepackage[T1]{fontenc}
\usepackage[utf8]{inputenc}
\usepackage{microtype}
\usepackage[colorlinks=true,urlcolor=blue]{hyperref}

\hypersetup{pdftitle={RankCloak Conceals the Surface Form of Synthetic Cryptographic Artifacts in Language Model Generated Text},pdfauthor={Alexander V. Mantzaris}}

\begin{document}

\begin{letter}{\\}

\opening{Dear Dr.\ Yang (editor),}

Please find the focused revision of \emph{RankCloak Conceals the Surface Form of Synthetic Cryptographic Artifacts in Language Model Generated Text} (submission 0170c7b2-1065-4344-a5d5-b839b51ac22c). Thank you and Reviewer 2 for the constructive second assessment. We retained the previous revisions and completed all the additional experiments requested.

RankCloak addresses an artifact-specific problem, transforming recognizable synthetic cryptographic artifacts into context-dependent token choices so their hexadecimal, UUID, base64, nonce, tag, ciphertext-like, or signature surface form is not directly exposed in generated cover text, while an authorized receiver reconstructs the artifact under a shared protocol configuration. Steganalytic classification is a distinct case study. The new results show that current detectors can often identify encoded covers, but the detection does not itself reveal or recover the concealed cryptographic artifacts.

The revision adds strict  deduplication with zero leakage under the declared audit, 250 human-authored secondary controls, genuine 'leave one model family out' tests, data-adaptive text detectors, an 'exact model aware' attacker, and validation frozen low FPR analysis with exact false positive counts. It also adds entropy gated embedding and a matched comparison of approximately 4-bit and 8-bit quantized versions of the same pinned Qwen model, including recovery and cross quantization detection. The complete declared computational sets and their validation finished successfully without generation execution failures.
The paper does not claim a cryptographic secrecy proof, universal undetectability, or robustness under every text transformation.

The existing GitHub and Zenodo statements remain with the final V3 Zenodo is saved and ready.

Thank you for your time.

Sincerely,

Alexander V. Mantzaris, Associate Professor \\
School of Data, Mathematical, and Statistical Sciences, UCF \\
4000 University Blvd, Orlando, FL 32817, USA \\
alexander.mantzaris@ucf.edu \\
TCII 211C, +1-407-692-0280
\end{letter}
\end{document}
